% Chapter 3: The Formal Universe
% Part I: Foundations

\chapter{형식 우주}
\label{ch:formal-universe}

\begin{quote}
\textit{``정의(definition)란 무엇인가를 이름 짓는 것이 아니라, 무엇을 원초적으로 받아들이는지를 선언하는 것이다.''}
\end{quote}

\section{형식 우주를 구성하는 원리}
\label{sec:design-principles}

수학적 이론을 세우는 첫 번째 단계는 이론이 다루는 세계를 형식적으로 규정하는 것이다. 이를 \textbf{형식 우주(formal universe)}라 부른다. 형식 우주의 설계에는 두 가지 상충하는 요구가 있다: 충분히 풍부해서 이론이 포착하고자 하는 현상을 표현할 수 있어야 하고, 동시에 충분히 절제해서 불필요한 구조를 전제하지 않아야 한다.

SCC의 형식 우주는 이 균형을 반영한다. 여섯 개의 구성 요소를 포함하며, 각각은 이론에서 대체 불가능한 고유 역할을 담당한다. 이전 장에서 직관적으로 소개한 응집 장, 연산자, 전달 핵은 이 장에서 그 정확한 수학적 지위를 부여받는다.

형식 우주의 설계에서 핵심적인 구별은 \textbf{원시 개념(primitives)}과 \textbf{파생 개념(derived notions)}의 구별이다. 원시 개념은 형식 우주에 직접 포함되어 이론의 출발점이 되는 구조이다. 파생 개념은 원시 개념으로부터 정의되어 구성되는 이차적 구조이다. 이 구별은 단순히 정리 정돈의 문제가 아니다---어떤 것을 원시 개념으로 받아들이는가는 이론의 존재론적 약속(ontological commitment)을 결정한다.

SCC의 형식 우주 설계는 다음 세 가지 원리에 의해 지배된다:

\begin{enumerate}
    \item \textbf{최소 원시성(Minimal Primitivity).} 원시 개념의 수를 최소화한다. 다른 원시 개념으로부터 구성할 수 있는 것은 파생 개념으로 처리한다. 이는 이론의 존재론적 부담을 경량화하고, 구성 요소 간 의존 구조를 명확히 한다.

    \item \textbf{기능적 정당성(Functional Justification).} 모든 원시 개념은 이론 내에서 대체 불가능한 기능적 역할---에너지, 술어(predicate), 또는 공리적 구조에의 참여---을 가져야 한다. 역할이 없는 구성 요소는 형식 우주에서 제거한다.

    \item \textbf{층 분리(Layer Separation).} 존재론(무엇이 있는가), 공리론(무엇이 참인가), 실현(어떻게 구체화하는가), 구현(어떻게 계산하는가)은 엄격히 분리된 계층이다. 형식 우주는 존재론 계층에 속한다.
\end{enumerate}

이 원리들의 적용은 이론 발전 과정에서 실제적인 결과를 낳았다. 두 개의 연산자---전이 연산자 $\mathbf{T}_t$와 공-소속 연산자 $\mathbf{C}_t$---가 원래 형식 우주에 포함되어 있었으나, 기능적 정당성 원리에 의해 파생 진단(derived diagnostic)으로 강등되었다. 이에 대해서는 \S\ref{sec:what-was-removed}에서 상세히 논의한다.


\section{형식 우주 튜플}
\label{sec:formal-universe-tuple}

SCC 이론의 형식 우주는 다음과 같은 구조화된 튜플이다:
\begin{equation}
    \mathfrak{C}^{\mathrm{soft}} = \Big( T,\ \{X_t\}_{t \in T},\ \{u_t\}_{t \in T},\ \{\mathrm{Cl}_t\}_{t \in T},\ \{\mathbf{N}_t, \mathbf{D}_t\}_{t \in T},\ \{\mathbf{M}_{t \to s}\}_{t,s \in T} \Big)
    \label{eq:formal-universe-ch3}
\end{equation}

이 튜플은 여섯 종류의 구성 요소---시간, 공간, 장, 닫힘, 인접과 구별, 시간적 전달---를 담고 있다. 각각을 순서대로 살펴보자.


\subsection{시간 지표 집합 $T$}

$T$는 \textbf{전순서(totally ordered)} 집합으로, 시간적 순서가 잘 정의되어 있다. 최소의 경우 $T$는 유한 이산 수열 $\{0, 1, 2, \ldots, n\}$이 될 수 있고, 일반적으로는 후속(succession) 개념이 부여된 임의의 전순서 집합이 될 수 있다.

이론은 연속 시간을 요구하지 않는다. 그러나 시간적 순서가 잘 정의되어 있을 것을 요구한다---이것은 시간적 전달 핵 $\mathbf{M}_{t \to s}$가 의미를 가지기 위한 최소 조건이다.


\subsection{관계적 지지 공간 $X_t$}

시각 $t \in T$ 마다 $X_t$는 \textbf{관계적 지지 공간(relational support space)}이다. 이것은 응집이 정의되는 지점(site)들의 집합이다. 시각 도메인에서 $X_t$는 공간 위치의 격자일 수 있고, 더 추상적인 도메인에서는 관계적 위치(relational loci)들의 임의의 유한 또는 가산 집합일 수 있다.

$X_t$는 시간에 따라 변할 수 있다. 이 일반성은 인지적 또는 지각적 시스템에 사용 가능한 관계적 지지가 시간에 따라 변할 수 있기 때문에 필요하다.

\begin{center}
\fbox{\parbox{0.85\textwidth}{\centering
\textbf{$X_t$의 존재론적 지위}\\[6pt]
$X_t$의 지점들은 전-주어진 객체가 아니라 \emph{관계적 위치}이다.\\
이론의 존재론적 원시자는 응집 장 $u_t$와 그 관계적 조직이지,\\
지점들의 개별적 정체성이 아니다.\\[3pt]
픽셀 격자의 이산 구조가 영상 분석을 전-주어진 시각 객체에\\
약속하지 않는 것과 같은 의미에서, $X_t$의 이산 구조는\\
이론을 전-주어진 객체에 약속하지 않는다.
}}
\end{center}


\subsection{부드러운 응집 장 $u_t$}

시각 $t \in T$ 마다
\begin{equation}
    u_t : X_t \to [0,1]
    \label{eq:cohesion-field}
\end{equation}
은 \textbf{부드러운 응집 장(soft cohesion field)}이다. 이것은 전-객체적 응집의 원초적 운반자이다. $u_t(x)$의 값은 지점 $x$가 시각 $t$에서 응집적 형성에 참여하는 정도를 나타낸다.

Chapter~1에서 강조했듯이, 응집 장은 다음의 어느 것도 아니다:
\begin{itemize}
    \item 사후 확률(posterior probability)이 아니다
    \item 클래스 소속 점수(class membership score)가 아니다
    \item 분할 마스크(segmentation mask)가 아니다
\end{itemize}
응집 장은 이론의 \textbf{일차적 존재론적 실체}이다: 닫힘, 구별, 경계, 지속성을 포함한 모든 추가 구조가 이 연속적 장으로부터 파생된다.


\subsection{부드러운 닫힘 연산자 $\mathrm{Cl}_t$}

시각 $t \in T$ 마다
\begin{equation}
    \mathrm{Cl}_t : [0,1]^{X_t} \to [0,1]^{X_t}
    \label{eq:closure-operator}
\end{equation}
는 \textbf{부드러운 닫힘 연산자(soft closure operator)}이다. 응집 장을 그 \emph{관계적으로 완성된 형태}로 매핑한다: 모든 지점의 응집도가 관계적 이웃으로부터 받는 지지를 반영하도록 갱신된 장이다.

닫힘은 응집이 \emph{자기-유지적}이 되는 메커니즘이다. 직관적으로, $\mathrm{Cl}_t(u)$는 ``$u$가 자기 자신의 관계 구조에 비추어 스스로를 어떻게 볼 것인가''에 대한 답이다.

닫힘 연산자의 형식적 성질은 공리 그룹 A(Chapter~\ref{ch:axioms})에서 규정된다. 이 장에서는 두 가지 핵심 사항만 강조해둔다:

\textbf{(1) 비멱등성(Non-idempotence).} 고전적 위상학적 닫힘은 멱등이다: $\mathrm{Cl}^2 = \mathrm{Cl}$. SCC의 닫힘은 의도적으로 멱등이 아니다. 대신 반복 적용이 수렴하는 \emph{안정화 경향}(stabilization tendency)을 가진다. 이는 전-객체 수준에서 닫힘이 완성된 사실이 아니라 진행 중인 과정이기 때문이다.

\textbf{(2) 자기-참조적 구조.} 닫힘 연산자는 장 $u$에 적용되어 갱신된 장 $\mathrm{Cl}(u)$를 산출하며, 이 갱신은 $u$ 자체에 의해 결정된 관계 구조를 통해 수행된다. 이 자기-참조적 구조는 SCC를 표준 위상 장 이론과 구별하는 핵심 특성이다.


\subsection{부드러운 인접 핵 $\mathbf{N}_t$}

시각 $t \in T$ 마다
\begin{equation}
    \mathbf{N}_t : X_t \times X_t \to [0, \infty)
    \label{eq:adjacency-kernel}
\end{equation}
는 \textbf{부드러운 인접 핵(soft adjacency kernel)}이다. $\mathbf{N}_t(x,y)$는 지점 $x$와 $y$ 사이의 관계적 근접성 또는 국소적 커플링의 정도를 측정한다.

인접 핵은 닫힘과 구별이 구축되는 \textbf{기질(substrate)}이다. 두 지점이 인접하다는 것은 그들이 서로의 응집 상태에 직접적으로 영향을 미칠 수 있다는 의미이다. 그러나 인접은 추이적(transitive)이 아니다: $x$가 $y$에 인접하고 $y$가 $z$에 인접하더라도, $x$가 $z$에 인접할 필요는 없다. 전역적 연결성(co-belonging)은 국소적 인접으로부터 구축되어야지, 인접의 원시적 성질로 전제되어서는 안 된다.


\subsection{부드러운 구별 연산자 $\mathbf{D}_t$}

시각 $t \in T$ 마다
\begin{equation}
    \mathbf{D}_t : X_t \times [0,1]^{X_t} \to [0,1]
    \label{eq:distinction-operator}
\end{equation}
는 \textbf{부드러운 구별 연산자(soft distinction operator)}이다. 지점 $x$와 참조 장(전형적으로 보수 장 $1 - u_t$, 즉 외부를 표현하는 장)이 주어졌을 때, $\mathbf{D}_t(x; 1-u_t)$는 지점 $x$가 외부에 대해 \emph{구조적으로 비대칭적}인 정도를 측정한다.

구별은 단순한 국소적 대비(local contrast)가 아니다. 이것은 응집 장의 내부가 비-응집 영역 전체에 대해 가지는 \textbf{구조적 비대칭}이다. 이 비대칭이 없으면, 응집 장은 구별 없는 확산 패턴이 되어 의미 있는 내부/외부 분절이 불가능해진다.

구별 연산자의 입력에 보수 장 $1-u_t$가 포함된다는 사실은 중요한 자기-참조적 구조를 부여한다: 장 $u_t$ 자체가 ``외부가 무엇인가''를 정의하며, 구별 연산자는 장이 그 자체가 정의한 외부에 대해 비대칭적인지를 측정한다.


\subsection{시간적 전달 핵 $\mathbf{M}_{t \to s}$}

시간 쌍 $t, s \in T$ 마다
\begin{equation}
    \mathbf{M}_{t \to s} : X_t \times X_s \to [0,1]
    \label{eq:transport-kernel}
\end{equation}
는 \textbf{시간적 전달(temporal transport) 또는 계승(inheritance) 핵}이다. $\mathbf{M}_{t \to s}(x,y)$는 시각 $t$에서 지점 $x$의 응집적 역할이 시각 $s$에서 지점 $y$에 의해 계승되는 정도를 나타낸다.

시간적 전달은 SCC의 시간적 정체성을 정의하는 메커니즘이다: \emph{시간을 통해 같은 형성}이란, 응집적 핵심이 전달 하에서 구조적으로 계승되는 형성이다. 이것은 추적 ID에 의한 점별 정체성(pointwise identity)과 근본적으로 다르며, 철학적 \textbf{발생적 정체성(genidentity)} 개념과 더 가깝다.

시간적 전달은 네 가지 일반성을 허용한다:
\begin{itemize}
    \item \textbf{일대다(one-to-many)}: 시각 $t$의 응집 영역이 시각 $s$의 여러 지점으로 분산될 수 있다
    \item \textbf{다대일(many-to-one)}: 시각 $t$의 여러 지점이 시각 $s$의 단일 계승자로 수렴할 수 있다
    \item \textbf{부분적(partial)}: 응집 내용의 일부가 계승자 없이 소실될 수 있다
    \item \textbf{확률적(probabilistic)}: 전달이 결정적일 필요가 없다
\end{itemize}

요구되는 것은 전달이 응집의 \emph{구조적 조직}을 보존하는 것이지, 개별 지점의 정체성을 보존하는 것이 아니다.


\subsection{여섯 구성 요소의 전체 그림}

형식 우주의 여섯 구성 요소를 그 역할과 함께 표~\ref{tab:formal-components}에 요약한다.

\begin{table}[ht]
\centering
\caption{형식 우주 $\mathfrak{C}^{\mathrm{soft}}$의 구성 요소}
\label{tab:formal-components}
\begin{tabular}{llll}
\hline
\textbf{구성 요소} & \textbf{유형} & \textbf{역할} & \textbf{관련 에너지/술어} \\
\hline
$T$ & 전순서 집합 & 시간적 순서 & $\mathcal{E}_{\mathrm{tr}}$ \\
$X_t$ & 유한/가산 집합 & 관계적 지지 & (기질) \\
$u_t : X_t \to [0,1]$ & 함수 & 응집 장 (원초적 존재자) & 모든 항 \\
$\mathrm{Cl}_t$ & 연산자 & 관계적 자기-완성 & $\mathcal{E}_{\mathrm{cl}}$, $\mathsf{Bind}$ \\
$\mathbf{N}_t$ & 핵(kernel) & 국소적 관계 구조 & $\mathcal{E}_{\mathrm{bd}}$ \\
$\mathbf{D}_t$ & 연산자 & 내부/외부 비대칭 & $\mathcal{E}_{\mathrm{sep}}$, $\mathsf{Sep}$ \\
$\mathbf{M}_{t \to s}$ & 핵(kernel) & 시간적 구조 계승 & $\mathcal{E}_{\mathrm{tr}}$, $\mathsf{Persist}$ \\
\hline
\end{tabular}
\end{table}

% Figure placeholder
\begin{figure}[t]
\centering
\fbox{\parbox{0.85\textwidth}{\centering\vspace{5cm}
\textbf{[그림 3.1]} 형식 우주 다이어그램.\\[6pt]
중앙: 응집 장 $u_t$ (원초적 존재자).\\
좌측: 닫힘 $\mathrm{Cl}_t$ (자기-완성), 인접 $\mathbf{N}_t$ (국소 구조).\\
우측: 구별 $\mathbf{D}_t$ (내외부 비대칭).\\
하단: 전달 $\mathbf{M}_{t \to s}$ (시간적 계승).\\
화살표: 의존 관계 ($\mathrm{Cl}_t$와 $\mathbf{D}_t$는 $\mathbf{N}_t$와 $u_t$에 의존).
\vspace{0.5cm}}}
\caption{형식 우주 $\mathfrak{C}^{\mathrm{soft}}$의 구조. 응집 장 $u_t$는 중앙의 원초적 존재자이며, 모든 연산자와 핵은 $u_t$와 상호작용한다. 닫힘과 구별은 인접 구조 위에 구축되고, 시간적 전달은 시간 쌍 사이를 연결한다.}
\label{fig:formal-universe-diagram}
\end{figure}


\subsection{예제: $5 \times 5$ 격자에서의 형식 우주}

구체적 예를 통해 형식 우주를 이해해보자. $5 \times 5$ 격자 그래프를 생각하자. 이 경우:

\begin{itemize}
    \item $T = \{0\}$ (단일 시점)
    \item $X_0 = \{(i,j) : 1 \leq i,j \leq 5\}$, 25개 지점
    \item $\mathbf{N}_0((i,j),(i',j')) = 1$ if $|i-i'| + |j-j'| = 1$, else $0$ (4-연결 인접)
\end{itemize}

예를 들어, 격자 중앙에 ``덩어리(bump)''가 있는 응집 장을 생각하자:
\begin{equation}
    u_0(i,j) = \exp\Big(-\frac{(i-3)^2 + (j-3)^2}{2}\Big)
    \label{eq:bump-example}
\end{equation}
이 장은 중앙 $(3,3)$에서 $u \approx 1$이고, 가장자리에서 $u \approx 0$이다. 중앙의 높은 값은 ``깊은 내부''를, 중간 값은 ``경계 대역''을, 낮은 값은 ``외부''를 나타낸다.

이 장에 닫힘 연산자를 적용하면 어떻게 될까? 닫힘은 각 지점의 응집도를 이웃의 평균에 기반하여 갱신한다. 중앙 지점 $(3,3)$은 4개의 이웃이 모두 높은 응집도를 가지므로 강한 관계적 지지를 받는다---닫힘 후에도 높은 값을 유지한다. 반면 모서리 지점 $(1,1)$은 이웃 중 높은 응집도를 가진 것이 없으므로 관계적 지지가 약하다---닫힘 후에도 낮은 값을 유지한다. 경계 지점은 내부와 외부 이웃을 함께 가지므로, 이들의 닫힘 값은 이 상충하는 지지를 반영하는 중간 값이 된다.

구별 연산자 $\mathbf{D}_0(x; 1-u_0)$는 이 장의 내부/외부 비대칭을 측정한다. 깊은 내부 지점 $(3,3)$에서 이웃의 응집도 평균 $P(u)(3,3)$은 높고 외부 장 평균 $P(1-u)(3,3)$은 낮으므로, 구별 값은 높다---이 지점은 외부에 대해 강하게 비대칭적이다. 외부 지점 $(1,1)$에서는 그 반대이다: 자기 자신의 응집도가 낮고 이웃도 낮으므로, 외부 장에 대한 비대칭이 약하다.

이 간단한 예제에서도 형식 우주의 구성 요소들이 협력하여 형성의 내부 구조---깊은 내부, 경계 대역, 외부의 분절---를 포착하는 것을 확인할 수 있다.


\subsection{수치 계산: $5 \times 5$ 격자의 연산자}
\label{sec:5x5-numerical}

위의 Gaussian 덩어리 예제를 수치적으로 전개해보자. 이것은 형식 우주의 구성 요소들이 구체적 수치 수준에서 어떻게 작동하는지를 보여주는 완전한 워크스루(walkthrough)이다.

\textbf{단계 1: 응집 장의 구체적 값.} 식~\eqref{eq:bump-example}에서 중앙 $(3,3)$을 기준으로 주요 지점의 값을 계산하면:
\begin{center}
\begin{tabular}{lcc}
\hline
\textbf{지점} & \textbf{$u_0(i,j)$} & \textbf{해석} \\
\hline
$(3,3)$ & $e^0 = 1.000$ & 깊은 내부 (핵심) \\
$(3,4)$, $(4,3)$ & $e^{-1/2} \approx 0.607$ & 내부 \\
$(4,4)$, $(2,2)$ & $e^{-1} \approx 0.368$ & 경계 대역 \\
$(3,5)$, $(5,3)$ & $e^{-2} \approx 0.135$ & 외부 경계 \\
$(1,1)$, $(5,5)$ & $e^{-4} \approx 0.018$ & 외부 \\
\hline
\end{tabular}
\end{center}

\textbf{단계 2: 닫힘 연산자의 적용.} 닫힘의 잠정적 실현(provisional realization, Part~II에서 정식으로 도입)은 다음 형태이다:
\begin{equation}
    \mathrm{Cl}_t(u)(x) = \sigma\!\Big( a_{\mathrm{cl}} \big( (1-\eta) \cdot u(x) + \eta \cdot (Pu)(x) - \tau \big) \Big)
    \label{eq:cl-provisional-ch3}
\end{equation}
여기서 $\sigma(z) = 1/(1+e^{-z})$는 로지스틱 시그모이드, $P$는 행-정규화된 인접 집성(aggregation) 연산자, $a_{\mathrm{cl}} = 3$, $\eta = 0.5$, $\tau = 0.5$로 놓자.

중앙 지점 $(3,3)$에서 $(Pu)(3,3)$는 네 이웃 $\{(2,3),(4,3),(3,2),(3,4)\}$의 평균이므로:
\[
    (Pu)(3,3) = \frac{u(2,3) + u(4,3) + u(3,2) + u(3,4)}{4} = \frac{4 \times 0.607}{4} = 0.607
\]
따라서:
\[
    \mathrm{Cl}(u)(3,3) = \sigma\big(3 \cdot (0.5 \times 1.000 + 0.5 \times 0.607 - 0.5)\big) = \sigma(3 \times 0.304) = \sigma(0.911) \approx 0.713
\]

외부 지점 $(1,1)$에서는 이웃이 $(1,2)$와 $(2,1)$뿐이고(모서리), 두 이웃 모두 $u \approx 0.135$이므로:
\[
    (Pu)(1,1) \approx 0.135, \qquad \mathrm{Cl}(u)(1,1) = \sigma\big(3(0.5 \times 0.018 + 0.5 \times 0.135 - 0.5)\big) = \sigma(-1.27) \approx 0.219
\]

이 결과는 닫힘의 핵심 성질을 보여준다: 관계적 지지가 강한 지점(중앙)은 높은 닫힘 값을 얻고, 약한 지점(모서리)은 낮은 닫힘 값을 얻는다.

\textbf{단계 3: 구별 연산자의 적용.} 구별의 잠정적 실현은:
\begin{equation}
    \mathbf{D}_t(x; 1-u) = \sigma\!\Big( a_D \big( (Pu)(x) - \lambda_D \cdot (P(1-u))(x) \big) - \tau_D \Big)
    \label{eq:d-provisional-ch3}
\end{equation}
$a_D = 3$, $\lambda_D = 1.5$, $\tau_D = 0$으로 놓자. 중앙 $(3,3)$에서:
\[
    (P(1-u))(3,3) = 1 - 0.607 = 0.393
\]
\[
    \mathbf{D}(3,3; 1-u) = \sigma\big(3(0.607 - 1.5 \times 0.393)\big) = \sigma(3 \times (-0.0825)) = \sigma(-0.248) \approx 0.438
\]

흥미롭게도, 중앙 지점에서의 구별은 ``중간'' 수준이다. 이는 중앙이 깊은 내부이지만 이웃 모두가 높은 응집도를 가져 \emph{외부와의 비대칭이 국소적으로는 약하기 때문}이다. 더 큰 격자에서 핵심 지점은 구별 값이 더 높아지는데, 이는 경계에서 급격한 응집 변화가 구별에 기여하기 때문이다.

\textbf{단계 4: 닫힘 에너지의 계산.} 닫힘 에너지 $\mathcal{E}_{\mathrm{cl}} = \sum_x (u(x) - \mathrm{Cl}(u)(x))^2$를 주요 지점에서 계산하면:
\begin{center}
\begin{tabular}{lccc}
\hline
\textbf{지점} & $u(x)$ & $\mathrm{Cl}(u)(x)$ & $(u - \mathrm{Cl})^2$ \\
\hline
$(3,3)$ & 1.000 & 0.713 & 0.082 \\
$(3,4)$ & 0.607 & $\sim 0.594$ & 0.0002 \\
$(1,1)$ & 0.018 & 0.219 & 0.040 \\
\hline
\end{tabular}
\end{center}

외부 지점 $(1,1)$에서 $u < \mathrm{Cl}(u)$인 것은 시그모이드의 비선형성 때문이다---이는 닫힘이 외부 지점의 응집도를 약간 ``끌어올리는'' 효과를 가짐을 의미한다. 에너지 최소화 과정에서 장은 이러한 불일치를 줄이는 방향으로 진화한다.

% Figure placeholder
\begin{figure}[t]
\centering
\fbox{\parbox{0.85\textwidth}{\centering\vspace{6cm}
\textbf{[그림 3.4]} $5 \times 5$ 격자 수치 예제.\\[6pt]
좌상: 응집 장 $u_0$ 히트맵 (Gaussian bump).\\
우상: 닫힘 연산자 적용 후 $\mathrm{Cl}(u)$ 히트맵.\\
좌하: 구별 연산자 $\mathbf{D}(x; 1-u)$ 히트맵.\\
우하: 닫힘 에너지 $(u - \mathrm{Cl}(u))^2$의 지점별 분포.\\[3pt]
색상 스케일: 파랑(0)에서 빨강(1).\\
중앙 영역의 높은 응집, 균일한 닫힘, 경계에서의 높은 구별에 주목.
\vspace{0.5cm}}}
\caption{$5 \times 5$ 격자에서의 형식 우주 구성 요소 수치 계산. 동일한 Gaussian bump 응집 장에 대해 닫힘, 구별, 에너지를 시각화한다. 닫힘은 관계적 지지를 반영하고, 구별은 내외부 비대칭을 포착하며, 에너지는 자기-지지와의 편차를 측정한다.}
\label{fig:5x5-numerical}
\end{figure}

이 수치 예제의 핵심 교훈은 세 가지이다:
\begin{enumerate}
    \item 형식 우주의 모든 구성 요소는 구체적 수치로 계산 가능하다---추상적 존재에 그치지 않는다.
    \item 닫힘과 구별은 장의 서로 다른 구조적 측면을 포착한다: 닫힘은 자기-지지를, 구별은 내외부 비대칭을.
    \item 작은 격자에서도 형성의 내부 구조(핵심, 경계, 외부)가 수치적으로 명확히 드러난다.
\end{enumerate}


\section{원시 개념과 파생 개념}
\label{sec:primitive-vs-derived}

형식 우주에 포함된 여섯 가지 구성 요소는 이론의 \textbf{원시 개념}이다. 이들은 정의를 통해 구성되는 것이 아니라, 이론의 출발점으로 받아들여지는 기본 구조이다. 이 외의 모든 개념---핵심(Core), 경계 대역(Boundary Band), 공-소속(Co-belonging), 형태학적 질 측도($\mathcal{Q}_{\mathrm{morph}}$) 등---은 원시 개념으로부터 \emph{파생}된다.

표~\ref{tab:primitive-derived}는 SCC의 주요 개념들을 원시/파생으로 분류하는 완전한 레지스트리이다.

\begin{table}[ht]
\centering
\caption{SCC 이론의 원시 개념 및 파생 개념 레지스트리}
\label{tab:primitive-derived}
\begin{tabular}{lllp{5.5cm}}
\hline
\textbf{개념} & \textbf{지위} & \textbf{기호} & \textbf{정의/역할} \\
\hline
시간 지표 집합 & 원시 & $T$ & 전순서 집합 \\
관계적 지지 & 원시 & $X_t$ & 응집이 정의되는 지점들 \\
응집 장 & 원시 & $u_t$ & 응집적 참여 강도 $[0,1]$ \\
닫힘 연산자 & 원시 & $\mathrm{Cl}_t$ & 관계적 자기-완성 \\
인접 핵 & 원시 & $\mathbf{N}_t$ & 국소적 관계 구조 \\
구별 연산자 & 원시 & $\mathbf{D}_t$ & 내부/외부 비대칭 \\
전달 핵 & 원시 & $\mathbf{M}_{t \to s}$ & 시간적 구조 계승 \\
\hline
핵심 & 파생 & $\mathrm{Core}_t$ & $\{x : u_t(x) \geq \theta_{\mathrm{core}}\}$ \\
내부 & 파생 & $\mathrm{Int}_t$ & $\{x : u_t(x) \geq \theta_{\mathrm{in}}\}$ \\
경계 대역 & 파생 & $\mathrm{Bd}_t$ & $\{x : \theta_1 < u_t(x) < \theta_2\}$ \\
외부 & 파생 & $\mathrm{Ext}_t$ & $\{x : u_t(x) \leq \theta_{\mathrm{ext}}\}$ \\
기울기 지표 & 파생 & $g_t$ & $\sum_y \mathbf{N}_t(x,y)|u_t(x) - u_t(y)|$ \\
공-소속 & 파생 (진단) & $\mathbf{C}_t$ & 비국소적 구조 통합 측정 \\
전이 연산자 & 파생 (진단) & $\mathbf{T}_t$ & 경계 형태학 진단 \\
형태학적 질 & 파생 & $\mathcal{Q}_{\mathrm{morph}}$ & $H_0$ 지속성 기반 형태 평가 \\
진단 벡터 & 파생 & $\mathbf{d}$ & $(\mathsf{Bind}, \mathsf{Sep}, \mathsf{Inside}, \mathsf{Persist})$ \\
\hline
\end{tabular}
\end{table}

원시 개념과 파생 개념의 구별이 중요한 이유는 세 가지이다.

첫째, \textbf{존재론적 명료성}이다. 이론이 무엇을 기본적으로 전제하는지를 명시함으로써, 이론의 존재론적 약속을 투명하게 만든다. SCC가 전제하는 것은 시간, 공간, 연속 장, 닫힘, 인접, 구별, 전달이다. 핵심이나 경계 같은 파생 개념은 이 원시 개념들로부터 구성되므로, 추가적 존재론적 부담을 지지 않는다.

둘째, \textbf{이론적 경제성}이다. 파생 개념은 원시 개념으로부터의 정의를 변경함으로써 대안적 형태를 탐색할 수 있다. 예를 들어, 핵심의 임계값 $\theta_{\mathrm{core}}$를 바꾸어도 이론의 기본 구조는 변하지 않는다. 반면 원시 개념을 변경하면 이론 자체가 변한다.

셋째, \textbf{설계 결정의 추적 가능성}이다. 어떤 개념이 파생 개념의 지위에서 원시 개념으로 승격되거나, 그 반대로 강등되는 것은 이론에 대한 중요한 결정이며, 그 근거가 명확히 기록되어야 한다. SCC에서는 두 번의 강등이 있었으며(\S\ref{sec:what-was-removed}), 각각은 기능적 정당성 원리에 의해 정당화되었다.


\section{연산자 삼각형}
\label{sec:operator-triad}

SCC의 형식 우주에는 세 개의 연산자적 구조가 포함되어 있다: 닫힘 $\mathrm{Cl}_t$, 구별 $\mathbf{D}_t$, 인접 $\mathbf{N}_t$. 이들은 원래 세 가지 자기-참조적 역할을 담당하도록 설계되었으며, 이 역할 구조를 \textbf{연산자 삼각형(operator triad)}이라 부른다.

\begin{enumerate}
    \item \textbf{자기-완성(Self-completion)}: $\mathrm{Cl}_t$\\
    응집 장이 자기 자신의 관계 구조 하에서 스스로를 \emph{완성}하는 과정. ``내가 나의 이웃들로부터 어떤 지지를 받는가?''에 대한 답. 에너지 항 $\mathcal{E}_{\mathrm{cl}}$과 진단 요소 $\mathsf{Bind}$에 진입한다.

    \item \textbf{자기-대비(Self-contrast)}: $\mathbf{D}_t$ vs. $1-u_t$\\
    응집 장이 자기 자신이 정의한 외부에 대해 스스로를 \emph{대비}하는 과정. ``나는 나의 외부에 비해 구별되는가?''에 대한 답. 에너지 항 $\mathcal{E}_{\mathrm{sep}}$과 진단 요소 $\mathsf{Sep}$에 진입한다.

    \item \textbf{자기-통합(Self-integration)}: $\mathbf{C}_t$ (진단)\\
    응집 장이 자기 자신의 관계 구조를 통해 비국소적으로 통합되는 정도. ``서로 떨어진 두 지점이 같은 형성에 속하는가?''에 대한 답. 현재 이론에서는 에너지나 술어에 진입하지 않으며, 순수한 진단 역할만 담당한다.
\end{enumerate}

% Figure placeholder
\begin{figure}[t]
\centering
\fbox{\parbox{0.85\textwidth}{\centering\vspace{5cm}
\textbf{[그림 3.2]} 연산자 삼각형(Operator Triad).\\[6pt]
삼각형의 세 꼭짓점: 자기-완성($\mathrm{Cl}_t$), 자기-대비($\mathbf{D}_t$), 자기-통합($\mathbf{C}_t$).\\
좌측 변: $\mathrm{Cl}_t$와 $\mathbf{D}_t$는 에너지에 직접 진입 (변분적).\\
우측 변: $\mathbf{C}_t$는 진단 전용 (비변분적).\\
$\mathbf{N}_t$ (인접)는 삼각형의 중앙에 위치---세 연산자의 공통 기질.\\[3pt]
실선: 에너지/술어에 진입하는 연산자.\\
점선: 진단 역할만 하는 연산자.
\vspace{0.5cm}}}
\caption{연산자 삼각형. 닫힘(자기-완성)과 구별(자기-대비)은 에너지 범함수에 직접 진입하는 ``변분적'' 연산자이다. 공-소속(자기-통합)은 형성의 비국소적 구조를 분석하는 ``진단적'' 연산자이다. 인접은 세 연산자가 공유하는 관계적 기질이다.}
\label{fig:operator-triad}
\end{figure}


\subsection{왜 세 개인가?}

세 개의 자기-참조적 역할은 형성이 가질 수 있는 세 가지 독립적인 구조적 성질에 대응한다:

\textbf{닫힘}은 형성의 \emph{내부적 일관성}을 다룬다. 형성의 각 부분이 나머지 부분에 의해 지지되는가? 자기-지지가 없는 형성은 무작위적 패턴일 뿐이다.

\textbf{구별}은 형성의 \emph{외부적 분리}를 다룬다. 형성이 그 배경으로부터 구조적으로 분리되는가? 구별이 없는 형성은 배경에 녹아든 확산 패턴이다.

\textbf{통합}은 형성의 \emph{전역적 연결성}을 다룬다. 형성의 멀리 떨어진 부분들이 하나의 통일된 구조에 속하는가? 통합이 없는 형성은 우연히 가까이 있는 독립적 조각들의 집합일 수 있다.

이 세 역할은 개념적으로 독립이다. 자기-지지는 있으나 외부와 구별이 없는 장(균일한 상수 장 $u \equiv c$), 외부와 구별은 되나 내부적으로 일관되지 않는 장(무작위 패턴), 국소적으로 일관되고 구별되나 전역적으로 통합되지 않는 장(떨어진 두 덩어리)은 모두 가능한 상태이다.

네 번째 독립적 역할을 추가하면 이론은 \emph{초과결정(overdetermined)}될 위험이 있다. 세 가지 역할은 형성의 구조적 성질을 포괄적이면서도 경제적으로 포착한다. 실제로, 이론의 12차례 반복적 발전(iteration)에서 네 번째 연산자를 추가할 필요가 발견되지 않았다.


\subsection{변분적 연산자와 진단적 연산자의 구별}

연산자 삼각형에서 주목할 중요한 구별은 \textbf{변분적(variational)} 연산자와 \textbf{진단적(diagnostic)} 연산자의 구별이다.

변분적 연산자는 에너지 범함수에 직접 진입하여 형성의 최적화 과정에 참여한다. SCC에서 닫힘 $\mathrm{Cl}_t$와 구별 $\mathbf{D}_t$가 이에 해당한다:
\begin{align}
    \mathcal{E}_{\mathrm{cl}} &= \sum_x (u_t(x) - \mathrm{Cl}_t(u_t)(x))^2 & &\text{(닫힘이 에너지에 진입)} \label{eq:ecl-uses-cl} \\
    \mathcal{E}_{\mathrm{sep}} &= \sum_x u_t(x)(1 - \mathbf{D}_t(x; 1-u_t)) & &\text{(구별이 에너지에 진입)} \label{eq:esep-uses-d}
\end{align}

진단적 연산자는 에너지에 진입하지 않으며, 형성의 구조를 사후적으로 분석하는 데만 사용된다. 공-소속 $\mathbf{C}_t$가 이에 해당한다. $\mathbf{C}_t$는 형성 내부의 비국소적 연결 구조를 3차 이상의 크기 차이로 식별할 수 있지만, 그 정보는 에너지 최소화 과정에는 사용되지 않는다.

이 구별이 중요한 이유는 계산적이기도 하다. 변분적 연산자의 야코비안(Jacobian)은 에너지 기울기 계산에 필요하므로, 효율적으로 계산 가능해야 한다. 공-소속의 레졸벤트(resolvent) 실현은 야코비안 노름이 $\sim 9300$에 달하여, 만약 에너지에 진입한다면 수치적 불안정성의 원인이 될 수 있었다. 진단 역할로의 강등은 이 문제를 원천적으로 제거한다.


\section{시간적 전달: 추적 없는 항등성}
\label{sec:temporal-transport}

형식 우주의 마지막 구성 요소인 시간적 전달 핵 $\mathbf{M}_{t \to s}$는 이론에서 특별한 위치를 차지한다. 이것은 SCC의 시간적 정체성 개념---추적 ID에 의존하지 않는 구조적 계승으로서의 정체성---을 실현하는 수학적 장치이다.


\subsection{왜 전달이 필요한가}

앞의 다섯 구성 요소($T$, $X_t$, $u_t$, $\mathrm{Cl}_t$, $\mathbf{N}_t$/$\mathbf{D}_t$)만으로는 단일 시점에서의 형성 구조만을 기술할 수 있다. 그러나 형성이 진정한 의미에서 ``무언가''이려면, 시간을 통해 \emph{지속}해야 한다. 하나의 시점에서만 존재하는 응집 패턴은 우연적 배치이지, 안정화된 형성이 아니다.

시간적 지속을 기술하는 가장 단순한 방법은 ``같은 지점이 같은 값을 유지한다''는 점별 항등성이다. 그러나 이것은 SCC의 맥락에서 부적절하다:
\begin{itemize}
    \item 형성은 공간적으로 이동할 수 있다 (같은 형성이 다른 위치에)
    \item 형성은 변형될 수 있다 (같은 형성이 다른 형태로)
    \item $X_t$ 자체가 변할 수 있다 (다른 지지 공간 위의 같은 형성)
    \item 경계는 이동할 수 있으나 핵심은 보존될 수 있다
\end{itemize}

필요한 것은 \emph{구조적 계승}의 개념이다: 형성의 핵심적 응집 조직이 시간을 통해 전달되었는가?


\subsection{전달 핵의 역할}

전달 핵 $\mathbf{M}_{t \to s}(x,y) \in [0,1]$는 이 구조적 계승을 부호화한다. 값이 높으면 지점 $x$의 응집적 역할이 지점 $y$에 의해 강하게 계승됨을 의미하고, 값이 낮으면 약하게 계승됨을 의미한다.

핵심적으로, 전달 핵은 \textbf{부분-확률적(sub-stochastic)}이다:
\begin{equation}
    \sum_{y \in X_s} \mathbf{M}_{t \to s}(x,y) \leq 1 \quad \text{for all } x \in X_t
    \label{eq:sub-stochastic}
\end{equation}
등호가 아닌 부등호를 사용하는 것은 부분적 소실(partial dissipation)을 허용하기 위함이다: 응집 내용의 일부는 계승자 없이 소실될 수 있다. 완전한 보존을 강제하는 등호는 일반적 경우에 너무 강하다.

이 부분-확률 구조는 최적 전달(optimal transport) 이론과 자연스럽게 연결된다. 전달 핵은 측도 $u_t \cdot \mu_{X_t}$와 $u_s \cdot \mu_{X_s}$ 사이의 부분-확률적 커플링을 정의하며, 이는 \textbf{비균형 최적 전달(unbalanced optimal transport)}---질량의 생성과 소멸이 허용되는 전달 문제---의 틀에 속한다.


\subsection{발생적 정체성(Genidentity)}

SCC의 시간적 전달 핵이 구현하는 정체성 개념은 점별 정체성(``같은 좌표가 같은 값'')이 아니라 \textbf{발생적 정체성(genidentity)}이다. 이 개념은 Hans Reichenbach가 도입한 것으로, ``시간적으로 연장된 동일성은 연속적 상태들 사이의 물리적 연결에 의해 구성된다''는 관점이다.

SCC에서 이를 구체화하면:

\begin{center}
\fbox{\parbox{0.85\textwidth}{\centering
\textbf{SCC의 정체성 원리}\\[6pt]
시간을 통해 \emph{같은 형성}이란,\\
응집적 핵심(core)이 전달 핵 하에서 구조적으로 계승되는 형성이다.\\[3pt]
핵심이 계승되지 않으면, 형성은 소멸한 것이다.\\
경계가 이동하고 형태가 변해도, 핵심이 계승되면 같은 형성이다.
}}
\end{center}

이 정체성 개념은 추적 ID 기반 접근과 세 가지 점에서 다르다:

\textbf{첫째, 비전체적(non-total)이다.} 형성이 소멸하는 것---전달 핵의 핵심 영역에서의 계승 합이 0에 가까워지는 것---은 이론 내에서 자연스럽다. 추적 ID에서는 ``ID가 사라진다''는 것이 시스템 바깥의 사건이다.

\textbf{둘째, 부분적(partial)이다.} 형성의 일부만 계승될 수 있다. 형성이 분열하면 두 계승자가 각각 원래 형성의 부분을 계승한다. 추적 ID에서는 ``같은 물건이 두 개가 되었다''는 모순이다.

\textbf{셋째, 구조적(structural)이다.} 정체성은 좌표의 동일성이 아니라 응집 조직의 동일성에 의해 결정된다. 형성이 공간적으로 이동해도, 핵심의 관계적 조직이 보존되면 같은 형성이다.

% Figure placeholder
\begin{figure}[t]
\centering
\fbox{\parbox{0.85\textwidth}{\centering\vspace{5cm}
\textbf{[그림 3.3]} 시간적 전달과 구조적 계승.\\[6pt]
상단: 시각 $t$의 형성 (핵심, 경계 대역, 외부).\\
하단: 시각 $s$의 형성 (위치 이동, 형태 변형).\\
화살표: 전달 핵 $\mathbf{M}_{t \to s}$에 의한 계승.\\
핵심 $\to$ 핵심: 강한 계승 (굵은 화살표).\\
경계 $\to$ 경계: 약한 계승 (가는 화살표).\\
일부 경계 지점: 계승 없음 (소실, 점선).
\vspace{0.5cm}}}
\caption{시간적 전달 핵에 의한 구조적 계승. 형성의 핵심은 강하게 계승되고, 경계는 부분적으로 계승되며, 일부 영역은 소실된다. 위치 이동과 형태 변형에도 불구하고 핵심이 계승되므로, SCC의 정체성 원리에 의해 이것은 ``같은 형성''이다.}
\label{fig:temporal-transport-schematic}
\end{figure}


\subsection{시간적으로 연장된 형식 우주}

형식 우주에 시간적 전달을 포함하면, SCC는 단일 시점의 이론이 아니라 \textbf{시간적으로 연장된} 이론이 된다. 시간적 창(temporal window) $W \subseteq T$ 위에서 정의된 형성 가족 $\mathbf{u} = (u_t)_{t \in W}$와 그 사이의 전달 핵 가족 $(\mathbf{M}_{t \to s})_{t, s \in W}$가 함께, 시간을 통해 지속하는 응집적 형성의 완전한 기술을 제공한다.

이 시간적으로 연장된 구조는 proto-cohesion 진단 벡터의 네 번째 요소인 $\mathsf{Persist}$의 기반이 된다. $\mathsf{Persist}$는 형성의 핵심이 시간 창 전체에 걸쳐 전달 핵을 통해 얼마나 잘 계승되었는지를 측정한다. 이에 대한 수학적 정의는 Part~II에서 정식으로 다룬다.


\section{파생 기하학적 개념}
\label{sec:derived-geometric}

원시 개념인 응집 장 $u_t$로부터, 기하학적 및 형태학적 의미를 가진 여러 파생 개념을 구성할 수 있다. 이들은 형성의 내부 구조를 분석하고 진단하는 데 사용된다.


\subsection{핵심(Core)}

형성의 \textbf{핵심}은 응집적 참여가 높은 임계값을 초과하는 지점들의 집합이다:
\begin{equation}
    \mathrm{Core}_t(u_t) = \{x \in X_t \mid u_t(x) \geq \theta_{\mathrm{core}}\}
    \label{eq:core-def}
\end{equation}
여기서 $\theta_{\mathrm{core}}$는 1에 가까운 매개변수이다. 핵심은 형성의 \emph{깊은 내부}---가장 강하고 안정적으로 응집 조직에 통합된 지점들---를 대표한다.

핵심은 시간적 지속성에서 특별한 역할을 한다. 형성의 정체성은 핵심의 계승에 의해 정의되기 때문에, 핵심을 잃은 형성은 지속을 중단한 것이다. 핵심은 형성의 \emph{구조적 핵(structural nucleus)}이다.


\subsection{내부(Interior)}

\textbf{내부}는 핵심보다 낮은 임계값으로 정의된다:
\begin{equation}
    \mathrm{Int}_t(u_t) = \{x \in X_t \mid u_t(x) \geq \theta_{\mathrm{in}}\}
    \label{eq:interior-def}
\end{equation}
여기서 $\theta_{\mathrm{in}} < \theta_{\mathrm{core}}$이다. 내부는 핵심을 포함하며, 강하지만 최대는 아닌 응집 참여를 보이는 주변 영역도 포함한다. 내부는 형성의 관계 구조가 실질적으로 자기-지지적인 영역이다.


\subsection{경계 대역(Boundary Band)}

\textbf{경계 대역}은 응집이 내부적 참여와 외부적 부재 사이의 중간인 영역이다:
\begin{equation}
    \mathrm{Bd}_t(u_t) = \{x \in X_t \mid \theta_1 < u_t(x) < \theta_2\}
    \label{eq:boundary-band-def}
\end{equation}
여기서 $0 < \theta_1 < \theta_2 < 1$이다.

경계 대역의 핵심적 특징은 이것이 날카로운 선이 아니라 \textbf{체적적 영역(volumetric region)}이라는 점이다. 위상학에서 경계는 전형적으로 여차원-1(codimension-one) 다양체이다---3차원에서 면, 2차원에서 곡선. 그러나 SCC에서 경계는 두께와 내부 구조를 가진 전이 대역이다. 두꺼운 경계 대역은 형성이 내부에서 외부로 점진적으로 전이함을 의미하고, 얇은 경계 대역은 급격한 전이를 의미한다. 두 경우 모두 정당하며, 이론은 어느 쪽도 선호하지 않는다.

경계 대역은 \textbf{기울기 지표(gradient indicator)}를 통해서도 특성화할 수 있다:
\begin{equation}
    g_t(x; u) = \sum_{y \in X_t} \mathbf{N}_t(x,y) \, |u_t(x) - u_t(y)|
    \label{eq:gradient-indicator}
\end{equation}
$g_t(x; u)$가 큰 지점은 응집의 급격한 변화가 일어나는 곳---형태학적 의미에서의 경계---이다.


\subsection{외부(Exterior)}

\textbf{외부}는 보수적으로 정의된다:
\begin{equation}
    \mathrm{Ext}_t(u_t) = \{x \in X_t \mid u_t(x) \leq \theta_{\mathrm{ext}}\}
    \label{eq:exterior-def}
\end{equation}
여기서 $\theta_{\mathrm{ext}}$는 0에 가까운 임계값이다. 외부는 형성의 구별이 측정되는 대상 도메인이다.


\subsection{형태학적 질 측도 $\mathcal{Q}_{\mathrm{morph}}$}

응집 장의 형태학적 질은 지속적 상동(persistent homology)에 기반한 측도로 평가된다:
\begin{equation}
    \mathcal{Q}_{\mathrm{morph}}(u_t) = \frac{\ell_{\max}(u_t) - c}{1 - c} \cdot \mathrm{Artic}(u_t)
    \label{eq:qmorph-def}
\end{equation}
여기서:
\begin{itemize}
    \item $c = \sum_x u_t(x) / n$은 체적 분율(volume fraction)
    \item $\ell_{\max}$는 초수준 집합 여과(superlevel-set filtration)의 $H_0$ 지속성 다이어그램에서 가장 긴 막대의 길이
    \item $\mathrm{Artic} = 1 - \ell_{\mathrm{second}} / \ell_{\max}$는 분절 비율---일차 형성이 이차 구조에 비해 얼마나 지배적인지를 측정
\end{itemize}

정규화 $(\ell_{\max} - c)/(1-c)$는 균일 장($u \equiv c$, 즉 $\ell_{\max} = c$)에서 $\mathcal{Q}_{\mathrm{morph}} = 0$이 되도록 보장한다. 곱 형태는 두 가지 독립적 요구를 포착한다: 정규화된 지속성은 일차 형성의 기선 위 견고성을 측정하고, 분절 비율은 단일 형성이 장을 얼마나 깨끗하게 지배하는지를 측정한다.

이 측도는 진단 벡터의 세 번째 요소 $\mathsf{Inside}$의 기반이 된다. 수학적 성질(QM1--QM4)은 Part~II에서 증명한다.


\section{제거된 것들과 그 이유}
\label{sec:what-was-removed}

SCC의 형식 우주는 이론 발전 과정에서 두 번의 주요 구조 변경을 겪었다. 두 경우 모두, 원래 형식 우주에 포함되어 있던 연산자가 기능적 정당성의 결여로 인해 파생 진단으로 강등되었다. 이 강등의 역사와 근거를 기록하는 것은, 형식 우주의 현재 형태가 임의적 선택이 아니라 반복적 검증의 결과임을 보이기 위함이다.


\subsection{전이 연산자 $\mathbf{T}_t$의 강등 (v2.0)}

\textbf{원래 역할.} 전이 연산자 $\mathbf{T}_t$는 이론의 첫 버전(v1.0)에서 형식 우주에 포함되어 있었으며, 경계 형태학---내부에서 외부로의 전이 구조---을 특성화하는 역할이 부여되었다. 두 개의 공리(T-Ax1, T-Ax2)가 제안되었다.

\textbf{강등 근거.} 이론의 다섯 차례 반복적 발전(I1--I5)을 거친 결과, $\mathbf{T}_t$는 다음과 같은 기능적 공백을 가지는 것이 확인되었다:

\begin{center}
\begin{tabular}{ll}
\hline
\textbf{기능적 역할} & \textbf{$\mathbf{T}_t$의 상태} \\
\hline
구체적 실현(realization) & 0개 \\
증명된 정리 & 0개 \\
에너지 항에의 진입 & 없음 \\
술어(predicate)에의 진입 & 없음 \\
\hline
\end{tabular}
\end{center}

다섯 차례의 반복에서 0개의 실현, 0개의 정리, 0개의 에너지 진입, 0개의 술어 진입은, 형식 우주에 포함될 기능적 정당성이 충족되지 않았음을 의미한다.

\textbf{대체 구조.} $\mathbf{T}_t$의 개념적 역할---경계 형태학 특성화---은 이미 다른 구조에 의해 충분히 수행되고 있었다:
\begin{itemize}
    \item 기울기 지표 $g_t(x; u)$: 경계에서의 응집 변화율
    \item 경계 대역 $\mathrm{Bd}_t$: 전이 영역의 직접적 정의
    \item 형태학적 질 측도 $\mathcal{Q}_{\mathrm{morph}}$: 내부 구조의 정량적 평가
\end{itemize}

$\mathbf{T}_t$는 형식 우주에서 제거되었으나, 개념적 역할은 보존되었다. 만약 미래에 경계 형태학에 특화된 변분적 역할---에너지 항에의 직접 진입---이 발견된다면, $\mathbf{T}_t$는 형식 우주로 복원될 수 있다. 현 시점에서 그러한 역할은 없다.


\subsection{공-소속 연산자 $\mathbf{C}_t$의 강등 (v2.1)}

\textbf{원래 역할.} 공-소속 연산자 $\mathbf{C}_t$는 비국소적 구조 통합---서로 떨어진 두 지점이 같은 형성에 속하는 정도---을 측정하는 것으로, v2.0까지 형식 우주에 포함되어 있었다. 공-소속이 형식 우주에 있었던 핵심 이유는 분리 술어(Sep)에 진입했기 때문이다.

\textbf{강등의 직접적 원인.} 이론의 8번째 반복(I8)에서, 분리 술어의 심각한 진단적 결함이 발견되었다. 원래의 $\mathbf{C}_t$-가중 분리 공식은:
\begin{equation}
    \mathsf{Sep}^{\text{old}} = \frac{\sum_x \mathbf{C}_t(x,x) \cdot \mathbf{D}_t(x; 1-u_t)}{\sum_x \mathbf{C}_t(x,x)}
    \label{eq:sep-old}
\end{equation}
이었다. 문제는 레졸벤트 실현 $\mathbf{C}_t = (I - \alpha W_{\mathrm{sym}})^{-1}$에서 대각 원소 $\mathbf{C}_t(x,x) \geq 1$이 \emph{모든} 지점에 대해 상당한 값을 가진다는 것이다---외부 지점에서도. 이로 인해 $\mathbf{C}_t$-가중 평균은 외부 지점에 과도한 가중치를 부여하여, 형성의 질과 무관하게 $\mathsf{Sep} \approx 0.5$를 산출했다. 이는 진단적으로 퇴화한 것이다.

\textbf{수정과 결과.} 분리 술어는 $u$-가중 형태로 교정되었다:
\begin{equation}
    \mathsf{Sep}^{\text{new}} = \frac{\sum_x u_t(x) \cdot \mathbf{D}_t(x; 1-u_t)}{\sum_x u_t(x)}
    \label{eq:sep-new}
\end{equation}
이 공식은 응집 값 $u_t(x)$를 자연스러운 중요도 가중치로 사용한다. 높은 응집도의 지점이 분리 점수에 더 많이 기여하고, 외부 지점($u \approx 0$)은 자동으로 무시된다.

이 교정의 결과, $\mathbf{C}_t$는 이론 내에서 유일한 술어 역할을 잃었다. $\mathbf{C}_t$는 원래 에너지에 진입하지 않도록 설계되어 있었으므로(계산적 효율을 위해), 술어 역할의 상실은 형식 우주에 포함될 기능적 정당성의 완전한 상실을 의미했다.

\textbf{진단적 가치의 보존.} $\mathbf{T}_t$와 마찬가지로, $\mathbf{C}_t$는 형식 우주에서 제거되었으나 개념적 역할과 공리(C1--C4)는 보존되었다. 레졸벤트 실현 $\mathbf{C}_t = (I - \alpha W_{\mathrm{sym}})^{-1}$은 형성 내부 지점 쌍과 경계 횡단 지점 쌍 사이에서 $10^3$ 이상의 크기 차이로 식별할 수 있다. 이 진단 능력은 파생 진단의 지위에서 온전히 사용 가능하다.

\textbf{교훈.} 두 강등은 공통된 교훈을 제공한다: \emph{개념적 역할은 형식 우주 포함의 필요조건이지, 충분조건이 아니다.} 충분조건은 기능적 역할---에너지 또는 술어에의 실질적 진입---이다. SCC의 형식 우주는 이 원리를 반복적으로 적용하여 정제된 결과이다.


\section{정리: 왜 이 여섯인가}
\label{sec:summary-why-six}

이 장의 논의를 종합하면, 형식 우주 $\mathfrak{C}^{\mathrm{soft}}$의 여섯 구성 요소 각각이 이론에 대체 불가능한 역할을 담당하는 이유를 간결히 정리할 수 있다:

\begin{enumerate}
    \item \textbf{$T$ (시간)}: 시간적 순서 없이는 지속성을 정의할 수 없다.
    \item \textbf{$X_t$ (공간)}: 응집이 정의될 지점들이 필요하다.
    \item \textbf{$u_t$ (응집 장)}: 이론의 원초적 존재자. 이것 없이는 이론이 없다.
    \item \textbf{$\mathrm{Cl}_t$ (닫힘)}: 자기-지지(self-support). 이것 없이는 응집이 무작위 패턴과 구별되지 않는다. 에너지 $\mathcal{E}_{\mathrm{cl}}$에 진입.
    \item \textbf{$\mathbf{N}_t / \mathbf{D}_t$ (인접/구별)}: 국소 구조와 내외부 분리. 이것 없이는 형성이 배경에서 분리되지 않는다. 에너지 $\mathcal{E}_{\mathrm{bd}}$, $\mathcal{E}_{\mathrm{sep}}$에 진입.
    \item \textbf{$\mathbf{M}_{t \to s}$ (전달)}: 시간적 계승. 이것 없이는 형성이 시간을 통해 지속할 수 없다. 에너지 $\mathcal{E}_{\mathrm{tr}}$에 진입.
\end{enumerate}

다른 것은 추가할 수 있지만(공-소속, 전이 연산자, 형태학적 진단), 이 여섯보다 적게는 만들 수 없다. 여섯은 이론이 요구하는 \textbf{최소한}이다.

\bigskip

다음 장에서는 이 형식 우주의 구성 요소들이 만족해야 하는 공리적 성질---닫힘의 축약성, 인접의 비추이성, 구별의 외부 감응성, 전달의 부분-확률성---을 엄밀히 규정한다. 형식 우주가 ``무엇이 있는가''를 선언했다면, 공리 그룹은 ``무엇이 참인가''를 선언한다.
